\chapter{Differential Forms and Integration}\label{differential-forms}
\chapterstatus{complete}

\section{Introduction}\label{differential-forms:section-introduction}

Differential forms are the integrands of geometry. A $k$-form is a quantity that can be integrated over oriented $k$-dimensional submanifolds: a $1$-form over curves (the
work of a force field along a path), a $2$-form over surfaces (the flux of a field through a surface), an $n$-form over an $n$-manifold (a volume). Why forms, and not functions,
are the right integrands is explained by the change of variables formula. Under a change of coordinates $x = \tau(y)$, the integral $\int f\,dx^1 \cdots dx^n$ becomes
$\int(f \circ \tau)\,|\det D\tau|\,dy^1 \cdots dy^n$, with the \emph{absolute value} of the Jacobian determinant, while the $n$-form $f\,dx^1 \wedge \cdots \wedge dx^n$ becomes
$(f \circ \tau)\det D\tau\,dy^1 \wedge \cdots \wedge dy^n$, with the determinant itself. The two rules agree for orientation-preserving changes of coordinates, and this is why an
$n$-form can be integrated over an oriented manifold without any further choice, such as coordinates or a metric. (A metric allows one to integrate functions instead:
Chapter~\ref{riemannian-geometry}.)

The \emph{exterior derivative} $d$ turns $k$-forms into $(k + 1)$-forms. On $\RR^3$ it unifies the gradient, the curl and the divergence of vector calculus, it satisfies
$d \circ d = 0$ (which contains $\operatorname{curl}\operatorname{grad} = 0$ and $\operatorname{div}\operatorname{curl} = 0$), and it makes \emph{Stokes' theorem}
(Theorem~\ref{differential-forms:theorem-stokes}) true: the integral of $d\omega$ over a manifold with boundary is the integral of $\omega$ over the boundary. This one statement
contains the fundamental theorem of calculus and the theorems of Green, Gauss and Stokes. Everything in Hodge theory is written in this language: de Rham cohomology
(Chapter~\ref{de-rham}), harmonic forms (Chapter~\ref{hodge-theorem}), and the decomposition of forms on a complex manifold into types (Chapter~\ref{dolbeault}).

The chapter goes as follows.
\begin{enumerate}
\item Differential forms, their three descriptions (sections of a bundle, alternating maps on tangent vectors, coefficients in charts) and their behaviour under changes of
coordinates; computations of wedge products and pullbacks (Section~\ref{differential-forms:section-forms}).
\item The exterior derivative, constructed and characterised step by step; the interior product, the Lie derivative and Cartan's formula; closed and exact forms and line
integrals.
\item Orientations, in terms of forms, of atlases and of unit normal fields (Section~\ref{differential-forms:section-orientations}).
\item Integration of forms, manifolds with boundary and Stokes' theorem, with worked computations (Section~\ref{differential-forms:section-integration}).
\end{enumerate}
References: \cite{lee-smooth-manifolds}, \cite{warner-foundations}, \cite{spivak-calculus-manifolds}.

\noindent\textbf{Prerequisites.} Chapter~\ref{vector-bundles} (Vector Bundles) and Chapter~\ref{measure-theory} (Measure and Integration).

\section{Differential forms}\label{differential-forms:section-forms}

\begin{definition}\label{differential-forms:definition-forms}
A \emph{differential $k$-form} on a smooth manifold $M$ is a smooth section of $\Lambda^kT^*M$; the space of $k$-forms is $\Omega^k(M)$, and $\Omega^\bullet(M) = \bigoplus_k\Omega^k(M)$ with the wedge
product (Proposition~\ref{multilinear-algebra:proposition-exterior-properties}) is a graded-commutative algebra over $C^\infty(M)$. By
Proposition~\ref{differential-forms:proposition-forms-descriptions} below, a $k$-form is the same as a smoothly varying family of alternating $k$-linear maps
$(T_pM)^k \to \RR$, or as a family of coefficient functions in charts transforming in a prescribed way. For $f \colon N \to M$ smooth, the \emph{pullback} $f^*\omega$ is defined by
$(f^*\omega)_p(v_1, \ldots, v_k) = \omega_{f(p)}(d_pf\,v_1, \ldots, d_pf\,v_k)$; it is an algebra homomorphism and $(g \circ f)^* = f^*g^*$.
\end{definition}

In a chart $(U, (x^1, \ldots, x^n))$, $dx^i$ is the differential of the coordinate function $x^i$, so $dx^i(\partial/\partial x^j) = \partial x^i/\partial x^j = \delta_{ij}$: the
$dx^i|_p$ form the basis of $T_p^*M$ dual to the basis $\partial/\partial x^j|_p$ of $T_pM$. For a multi-index $I = (i_1 < \cdots < i_k)$ we write
$dx^I = dx^{i_1} \wedge \cdots \wedge dx^{i_k}$; for a sequence $(i_1, \ldots, i_k)$ that is not increasing, $dx^{i_1} \wedge \cdots \wedge dx^{i_k}$ is $0$ if two indices coincide
and $\pm dx^I$ otherwise, with the sign of the permutation that sorts the indices.

\begin{proposition}\label{differential-forms:proposition-forms-descriptions}
Let $M$ be a smooth $n$-manifold. For each $p$, the fibre $\Lambda^kT_p^*M$ is the space of alternating $k$-linear maps $(T_pM)^k \to \RR$
(Proposition~\ref{multilinear-algebra:proposition-dual-exterior} and Notation~\ref{multilinear-algebra:notation-forms}). For a family $\omega = (\omega_p)_{p \in M}$ of such maps the
following are equivalent:
\begin{enumerate}
\item $\omega$ is a smooth section of $\Lambda^kT^*M$, that is, a $k$-form;
\item for all smooth vector fields $X_1, \ldots, X_k$ on an open set $U$, the function $p \mapsto \omega_p(X_1(p), \ldots, X_k(p))$ is smooth on $U$;
\item in every chart of some atlas, the \emph{coefficients} $\omega_I(p) = \omega_p(\partial/\partial x^{i_1}|_p, \ldots, \partial/\partial x^{i_k}|_p)$ are smooth.
\end{enumerate}
In a chart $\omega = \sum_I\omega_I\,dx^I$, uniquely. If $(y^1, \ldots, y^n)$ is a second chart, then on the overlap
\[
dx^I = \sum_J\det\Bigl(\frac{\partial x^I}{\partial y^J}\Bigr)dy^J, \qquad \omega^{(y)}_J = \sum_I\omega^{(x)}_I\det\Bigl(\frac{\partial x^I}{\partial y^J}\Bigr),
\]
where $\partial x^I/\partial y^J$ is the $k \times k$ matrix $(\partial x^{i_a}/\partial y^{j_b})_{a,b}$ and the superscripts indicate the chart. For $k = n$ this reads
$\omega^{(y)} = \omega^{(x)}\det(\partial x/\partial y)$, with the Jacobian determinant of the transition map. Conversely, functions $\omega^{(x)}_I$ given on the charts of an atlas
and satisfying this rule on all overlaps come from a unique $k$-form.
\end{proposition}

\begin{proof}
\emph{The expression in a chart.} Since the $dx^i|_p$ are the basis dual to the $\partial/\partial x^i|_p$, the $dx^I|_p$ form a basis of $\Lambda^kT_p^*M$
(Proposition~\ref{multilinear-algebra:proposition-exterior-properties}(3)), and for increasing $I$ and $J$,
$dx^I(\partial/\partial x^{j_1}, \ldots, \partial/\partial x^{j_k}) = \det(dx^{i_a}(\partial/\partial x^{j_b}))_{a,b}$, which is $1$ if $I = J$ and $0$ otherwise
(Notation~\ref{multilinear-algebra:notation-forms}). Evaluating $\omega_p = \sum_Ic_I\,dx^I|_p$ on $(\partial/\partial x^{j_1}, \ldots, \partial/\partial x^{j_k})$ therefore gives
$c_J = \omega_J(p)$: the coefficients are the $\omega_J$, and they are unique.

(1) $\Leftrightarrow$ (3). The smooth structure of $\Lambda^kT^*M$ (Proposition~\ref{vector-bundles:proposition-operations}) is the one for which the frames $(dx^I)$ of charts are
local trivialisations; a section is smooth if and only if its coefficients in these frames are smooth. (3) $\Rightarrow$ (2). On a chart domain write
$X_b = \sum_jX_b^j\,\partial/\partial x^j$. By multilinearity $\omega(X_1, \ldots, X_k) = \sum_I\omega_I\,dx^I(X_1, \ldots, X_k)$, and
$dx^I(X_1, \ldots, X_k) = \det(X_b^{i_a})_{a,b}$ is a polynomial in the smooth functions $X_b^j$. (2) $\Rightarrow$ (3). Take the $X_b$ to be coordinate vector fields.

\emph{Change of chart.} By the first part applied to the chart $y$, the coefficient of $dy^J$ in $dx^I$ is
$dx^I(\partial/\partial y^{j_1}, \ldots, \partial/\partial y^{j_k}) = \det(dx^{i_a}(\partial/\partial y^{j_b})) = \det(\partial x^{i_a}/\partial y^{j_b})$, since
$dx^i(\partial/\partial y^j) = \partial x^i/\partial y^j$ is the derivative of the function $x^i$ in the direction $\partial/\partial y^j$. Substituting into
$\omega = \sum_I\omega^{(x)}_I\,dx^I$ and comparing coefficients gives the rule for $\omega^{(y)}_J$. For $k = n$ there is one multi-index, $(1, \ldots, n)$. Conversely, if
coefficients satisfying the rule are given on the charts of an atlas, define $\omega_p = \sum_I\omega^{(x)}_I(p)\,dx^I|_p$ using any chart at $p$; the rule says exactly
that this does not depend on the chart, and $\omega$ is smooth by (3).
\end{proof}

Two operations are used constantly: the wedge product of forms, computed pointwise with the rules of the exterior algebra
($dx^i \wedge dx^j = -dx^j \wedge dx^i$, so $dx^i \wedge dx^i = 0$), and the pullback, computed by substituting: for $f \colon N \to M$,
$f^*(\sum_I\omega_I\,dx^{i_1} \wedge \cdots \wedge dx^{i_k}) = \sum_I(\omega_I \circ f)\,d(x^{i_1} \circ f) \wedge \cdots \wedge d(x^{i_k} \circ f)$, since $f^*$ is an algebra homomorphism with
$f^*(dg) = d(g \circ f)$ for functions (the chain rule).

\begin{example}\label{differential-forms:example-form-computations}
\begin{enumerate}
\item \emph{A wedge product.} On $\RR^3$, for $\alpha = x\,dx + y\,dy$ and $\beta = dx + z\,dz$,
\[
\alpha \wedge \beta = x\,dx \wedge dx + xz\,dx \wedge dz + y\,dy \wedge dx + yz\,dy \wedge dz = -y\,dx \wedge dy + xz\,dx \wedge dz + yz\,dy \wedge dz .
\]
On the vectors $e_1, e_2$ it takes the value $\alpha(e_1)\beta(e_2) - \alpha(e_2)\beta(e_1) = x \cdot 0 - y \cdot 1 = -y$, in accordance with the coefficient of $dx \wedge dy$.
\item \emph{Pullback to a curve.} For the helix $\gamma(t) = (\cos t, \sin t, t)$ and $\omega = z\,dx + x\,dy$,
$\gamma^*\omega = t\,d(\cos t) + \cos t\,d(\sin t) = (\cos^2t - t\sin t)\,dt$.
\item \emph{Pullback to the sphere.} Let $x(\vartheta, \varphi) = (\sin\vartheta\cos\varphi, \sin\vartheta\sin\varphi, \cos\vartheta)$ and $\omega = x\,dy \wedge dz + y\,dz \wedge dx + z\,dx \wedge dy$. By
Proposition~\ref{differential-forms:proposition-forms-descriptions}, the pullback of $dy \wedge dz$ is the Jacobian minor $\partial(y, z)/\partial(\vartheta, \varphi)$ times
$d\vartheta \wedge d\varphi$, and similarly for the others:
\[
x^*(dy \wedge dz) = \sin^2\vartheta\cos\varphi\,d\vartheta \wedge d\varphi, \quad x^*(dz \wedge dx) = \sin^2\vartheta\sin\varphi\,d\vartheta \wedge d\varphi, \quad
x^*(dx \wedge dy) = \sin\vartheta\cos\vartheta\,d\vartheta \wedge d\varphi .
\]
So $x^*\omega = (\sin^3\vartheta\cos^2\varphi + \sin^3\vartheta\sin^2\varphi + \sin\vartheta\cos^2\vartheta)\,d\vartheta \wedge d\varphi = \sin\vartheta\,d\vartheta \wedge d\varphi$, the area form of the
unit sphere in spherical coordinates.
\item \emph{Spherical coordinates in space.} For $(r, \vartheta, \varphi) \mapsto r\,x(\vartheta, \varphi)$, the pullback of $dx \wedge dy \wedge dz$ is the Jacobian determinant times
$dr \wedge d\vartheta \wedge d\varphi$, namely $r^2\sin\vartheta\,dr \wedge d\vartheta \wedge d\varphi$ (Exercise~\ref{differential-forms:exercise-spherical-volume}).
\end{enumerate}
\end{example}

\begin{example}\label{differential-forms:example-polar-forms}
In the polar coordinates $(r, \theta)$ of Example~\ref{smooth-manifolds:example-polar}, with $x = r\cos\theta$ and $y = r\sin\theta$, the rule of
Proposition~\ref{differential-forms:proposition-forms-descriptions} for $k = 1$ gives $dx = \cos\theta\,dr - r\sin\theta\,d\theta$ and $dy = \sin\theta\,dr + r\cos\theta\,d\theta$.
Solving,
\[
dr = \cos\theta\,dx + \sin\theta\,dy = \frac{x\,dx + y\,dy}{r}, \qquad d\theta = \frac{-\sin\theta\,dx + \cos\theta\,dy}{r} = \frac{x\,dy - y\,dx}{x^2 + y^2},
\]
the second being the angle form of Counterexample~\ref{differential-forms:counterexample-closed-not-exact}. Compare with the tangent vectors
$\partial/\partial r = \cos\theta\,\partial/\partial x + \sin\theta\,\partial/\partial y$ and $\partial/\partial\theta = -r\sin\theta\,\partial/\partial x + r\cos\theta\,\partial/\partial y$.
The formula $dr = \cos\theta\,dx + \sin\theta\,dy$ found in many texts is about the covector $dr$, not the vector $\partial/\partial r$: vectors change by the Jacobian matrix
$P = \partial(x, y)/\partial(r, \theta)$ (its columns give $\partial/\partial r$ and $\partial/\partial\theta$), and covectors by its inverse (the rows of $P^{-1}$ give $dr$ and $d\theta$).
The coefficients agree for $r$ and not for $\theta$ because $P$ is a rotation followed by $\operatorname{diag}(1, r)$: $\partial/\partial r$ is a unit vector orthogonal to the
circles $r = \text{const}$, so the Euclidean metric identifies it with $dr$, while $\partial/\partial\theta$ has length $r$ and corresponds to $r^2\,d\theta$. For $2$-forms the rule
gives
\[
dx \wedge dy = \det\frac{\partial(x, y)}{\partial(r, \theta)}\,dr \wedge d\theta = r\,dr \wedge d\theta,
\]
which is the change of variables $\int f\,dx\,dy = \int f\,r\,dr\,d\theta$ of Example~\ref{measure-theory:example-gaussian}, in the language of forms
(Proposition~\ref{differential-forms:proposition-integral-well-defined}).
\end{example}

\begin{theorem}\label{differential-forms:theorem-exterior-derivative}
There is a unique family of linear maps $d \colon \Omega^k(M) \to \Omega^{k+1}(M)$, for all $M$ and $k$, such that:
$df$ is the differential of $f$ for $f \in \Omega^0(M) = C^\infty(M)$; $d(\alpha \wedge \beta) = d\alpha \wedge \beta + (-1)^{|\alpha|}\alpha \wedge d\beta$; and $d \circ d = 0$. It is local, given in a chart by
$d(\sum_I\omega_I\,dx^I) = \sum_I d\omega_I \wedge dx^I$, and natural: $f^*(d\omega) = d(f^*\omega)$. Moreover, for vector fields,
\[
d\omega(X_0, \ldots, X_k) = \sum_i(-1)^i X_i\left(\omega(X_0, \ldots, \widehat{X_i}, \ldots, X_k)\right) + \sum_{i<j}(-1)^{i+j}\omega([X_i, X_j], X_0, \ldots, \widehat{X_i}, \ldots, \widehat{X_j}, \ldots, X_k).
\]
\end{theorem}

\begin{proof}
We write $\partial_l = \partial/\partial x^l$ in a chart, and $J \setminus j_m$ for the multi-index $J = (j_0 < \cdots < j_k)$ with $j_m$ removed.

\emph{Step 1: the formula in a chart and its coefficients.} On a chart domain $U$ define
\[
d_U\Bigl(\sum_I\omega_I\,dx^I\Bigr) = \sum_Id\omega_I \wedge dx^I = \sum_I\sum_{l=1}^n\partial_l\omega_I\,dx^l \wedge dx^I.
\]
Then $d_U(f\,dx^{i_1} \wedge \cdots \wedge dx^{i_k}) = df \wedge dx^{i_1} \wedge \cdots \wedge dx^{i_k}$ also for sequences that are not increasing, since both sides vanish for repeated
indices and change by the same sign when the indices are sorted. We compute the coefficients of $d_U\omega$: for $J = (j_0 < \cdots < j_k)$,
\[
(d_U\omega)(\partial_{j_0}, \ldots, \partial_{j_k}) = \sum_{m=0}^k(-1)^m\,\partial_{j_m}\omega_{J \setminus j_m}.
\]
Indeed the form $dx^l \wedge dx^I$ vanishes on $(\partial_{j_0}, \ldots, \partial_{j_k})$ unless $l \notin I$ and $\{l\} \cup I = J$ as sets, that is, unless $l = j_m$ and
$I = J \setminus j_m$ for some $m$; and then $dx^{j_m} \wedge dx^{J \setminus j_m} = (-1)^m\,dx^J$, moving $dx^{j_m}$ past the $m$ forms $dx^{j_0}, \ldots, dx^{j_{m-1}}$, which takes
the value $(-1)^m$. Conversely, collecting the terms, $\sum_J\bigl(\sum_m(-1)^m\partial_{j_m}\omega_{J \setminus j_m}\bigr)dx^J = \sum_I\sum_l\partial_l\omega_I\,dx^l \wedge dx^I$.

\emph{Step 2: the three rules hold in a chart.} For a function, $d_Uf = \sum_l\partial_lf\,dx^l$ is the differential $df$, by the dual basis property. For the Leibniz rule it suffices by
linearity to take $\alpha = f\,dx^I$ with $|I| = k$ and $\beta = g\,dx^J$; then $\alpha \wedge \beta = fg\,dx^I \wedge dx^J$ and
\[
d_U(\alpha \wedge \beta) = (g\,df + f\,dg) \wedge dx^I \wedge dx^J = (df \wedge dx^I) \wedge (g\,dx^J) + (-1)^k(f\,dx^I) \wedge (dg \wedge dx^J),
\]
moving the $1$-form $dg$ past the $k$-form $dx^I$. Finally
$d_Ud_U(f\,dx^I) = \sum_{m,l}\partial_m\partial_lf\,dx^m \wedge dx^l \wedge dx^I = 0$: the terms $(m, l)$ and $(l, m)$ cancel because $\partial_m\partial_lf = \partial_l\partial_mf$
(Theorem~\ref{real-analysis:theorem-schwarz}) while $dx^m \wedge dx^l = -dx^l \wedge dx^m$, and the terms with $m = l$ vanish.

\emph{Step 3: uniqueness and locality.} Let $D$ be any family of linear maps with the three properties. If $\omega$ vanishes on an open set $U \ni p$, choose a bump function
$\chi$ with support in $U$ and $\chi(p) = 1$ (Lemma~\ref{smooth-manifolds:lemma-bump-manifold}); then $\chi\omega = 0$, so $0 = D(\chi\omega) = d\chi \wedge \omega + \chi\,D\omega$, and at
$p$ this gives $(D\omega)_p = 0$. So $D\omega$ near $p$ depends only on $\omega$ near $p$. Now let $p$ lie in a chart domain $U$, choose bump functions $\chi$ and $\tilde\chi$ with
support in $U$, with $\chi = 1$ near $p$ and $\tilde\chi = 1$ near the support of $\chi$
(Corollary~\ref{smooth-manifolds:corollary-partition-applications}(1), applied to a compact neighbourhood of $\operatorname{supp}\chi$ in $U$). Then $\chi\omega = \sum_I(\chi\omega_I)\,d(\tilde\chi x^{i_1}) \wedge \cdots \wedge d(\tilde\chi x^{i_k})$
is a form on $M$ built from global functions, and the Leibniz rule and $D \circ D = 0$ give $D(\chi\omega) = \sum_Id(\chi\omega_I) \wedge d(\tilde\chi x^{i_1}) \wedge \cdots \wedge d(\tilde\chi x^{i_k})$, since
$D(d(\tilde\chi x^i)) = DD(\tilde\chi x^i) = 0$. Near $p$ this is $\sum_Id\omega_I \wedge dx^I$, and $D\omega = D(\chi\omega)$ near $p$ by locality. So $D$ is given by the coordinate formula
in every chart, and is unique.

\emph{Step 4: existence, and what happens on the intersection of two charts.} Define $d\omega$ on each chart domain $U$ by $d_U$. On the intersection of two chart domains $U$ and $V$,
the operator $d_V$ satisfies the three rules (Step 2, on the manifold $U \cap V$), so, writing $\omega = \sum_I\omega_I\,dx^I$ in the coordinates of $U$ and using
$d_V(dx^i) = d_V(d_Vx^i) = 0$ and the Leibniz rule,
\[
d_V\Bigl(\sum_I\omega_I\,dx^I\Bigr) = \sum_Id_V\omega_I \wedge dx^I = \sum_Id\omega_I \wedge dx^I = d_U\Bigl(\sum_I\omega_I\,dx^I\Bigr).
\]
So the local definitions agree on overlaps and glue to an operator $d$ on $M$, which satisfies the three rules because they are local.

\emph{Step 5: naturality.} For a function $g$ on $M$, $f^*(dg) = d(g \circ f)$ by the chain rule. Locally $\omega = \sum_I\omega_I\,dx^{i_1} \wedge \cdots \wedge dx^{i_k}$, so
$f^*\omega = \sum_I(\omega_I \circ f)\,d(x^{i_1} \circ f) \wedge \cdots \wedge d(x^{i_k} \circ f)$, and by the Leibniz rule and $d \circ d = 0$ on $N$,
$d(f^*\omega) = \sum_Id(\omega_I \circ f) \wedge d(x^{i_1} \circ f) \wedge \cdots \wedge d(x^{i_k} \circ f) = f^*\bigl(\sum_Id\omega_I \wedge dx^I\bigr) = f^*(d\omega)$.

\emph{Step 6: the invariant formula.} Call the right-hand side $T(X_0, \ldots, X_k)$. It is local in the $X_i$. It is $C^\infty(M)$-linear in $X_0$: replacing $X_0$ by $fX_0$, the
term $i = 0$ becomes $f$ times itself; the terms $i \geq 1$ acquire the extra summands $(-1)^i(X_if)\,\omega(X_0, \ldots, \widehat{X_i}, \ldots, X_k)$; and the bracket terms with
$i = 0$ acquire, since $[fX_0, X_j] = f[X_0, X_j] - (X_jf)X_0$, the extra summands $-(-1)^j(X_jf)\,\omega(X_0, X_1, \ldots, \widehat{X_j}, \ldots, X_k)$, which cancel the previous ones.
A similar check shows that $T$ changes sign when two arguments are exchanged, so it is $C^\infty(M)$-linear in every argument, like $d\omega$. Two such expressions agree if they
agree on coordinate fields, and by alternation it suffices to take $\partial_{j_0}, \ldots, \partial_{j_k}$ with $j_0 < \cdots < j_k$. For coordinate fields all brackets vanish,
and $T(\partial_{j_0}, \ldots, \partial_{j_k}) = \sum_m(-1)^m\partial_{j_m}\omega_{J \setminus j_m}$, which is $d\omega(\partial_{j_0}, \ldots, \partial_{j_k})$ by Step 1.
\end{proof}

\begin{example}\label{differential-forms:example-exterior-derivative-low}
For $k = 1$ the invariant formula reads $d\omega(X, Y) = X\omega(Y) - Y\omega(X) - \omega([X, Y])$, and in coordinates on $\RR^2$, for $\omega = P\,dx + Q\,dy$,
$d\omega = (\partial_xQ - \partial_yP)\,dx \wedge dy$: Step 1 with $J = (1, 2)$ gives the coefficient $\partial_1\omega_2 - \partial_2\omega_1$. This is the integrand of Green's theorem,
which is Stokes' theorem (Theorem~\ref{differential-forms:theorem-stokes}) for a domain in the plane.
\end{example}

\begin{definition}\label{differential-forms:definition-interior-lie}
For $X \in \mathfrak{X}(M)$ the \emph{interior product} $\iota_X \colon \Omega^k \to \Omega^{k-1}$ is $(\iota_X\omega)(X_1, \ldots, X_{k-1}) = \omega(X, X_1, \ldots, X_{k-1})$
(Definition~\ref{multilinear-algebra:definition-interior-product}), and the \emph{Lie derivative} is $L_X\omega = \frac{d}{dt}\big|_{t=0}\Phi_t^*\omega$, where $\Phi$ is the flow of $X$.
\end{definition}

The interior product inserts $X$ as the first argument; on a function it is $0$, and on a $1$-form $\iota_X\alpha = \alpha(X)$. By
Proposition~\ref{multilinear-algebra:proposition-interior-antiderivation} it is an antiderivation: $\iota_X(\alpha \wedge \beta) = \iota_X\alpha \wedge \beta + (-1)^k\alpha \wedge \iota_X\beta$ for a
$k$-form $\alpha$. For instance $\iota_{\partial_x}(dx \wedge dy) = dy$ and $\iota_{\partial_y}(dx \wedge dy) = -dx$.

The Lie derivative measures how $\omega$ changes along the flow of $X$. The flow is in general only defined for small $|t|$ near each point
(Theorem~\ref{smooth-manifolds:theorem-flow}), but this suffices: $(\Phi_t^*\omega)_p$ involves only the values of $\omega$ near $\Phi_t(p)$, it is defined for $|t|$ small, and it depends
smoothly on $(t, p)$, as its coefficients in a chart are built from the coefficients of $\omega$ and the derivatives of the flow. So $L_X\omega$ is again a $k$-form. On a function,
$L_Xf = \frac{d}{dt}f(\Phi_t(p))|_{t=0} = Xf$.

\begin{example}\label{differential-forms:example-lie-derivative}
On $\RR^2$ let $X = x\,\partial_x + y\,\partial_y$, whose flow is $\Phi_t(p) = e^tp$. Then $\Phi_t^*(dx \wedge dy) = d(e^tx) \wedge d(e^ty) = e^{2t}dx \wedge dy$, so
$L_X(dx \wedge dy) = 2\,dx \wedge dy$: the flow expands areas at the rate $2$. For the rotation field $Y = -y\,\partial_x + x\,\partial_y$, the flow consists of rotations, which preserve
$dx \wedge dy$, so $L_Y(dx \wedge dy) = 0$. Cartan's formula below gives the same results by pure algebra: $\iota_X(dx \wedge dy) = x\,dy - y\,dx$, whose differential is $2\,dx \wedge dy$, and
$\iota_Y(dx \wedge dy) = -y\,dy - x\,dx = -d(\frac12(x^2 + y^2))$, whose differential is $0$.
\end{example}

\begin{proposition}[Cartan's formulas]\label{differential-forms:proposition-cartan}
$L_X = d\iota_X + \iota_Xd$ on forms; $L_X$ commutes with $d$; $\frac{d}{dt}\Phi_t^*\omega = \Phi_t^*L_X\omega$; and $L_{[X,Y]} = L_XL_Y - L_YL_X$.
\end{proposition}

\begin{proof}
\emph{Step 1: $L_X$ is a local derivation of degree $0$ commuting with $d$.} Since $\Phi_t^*(\alpha \wedge \beta) = \Phi_t^*\alpha \wedge \Phi_t^*\beta$, differentiating at $t = 0$ (the product
rule, applied pointwise to the coefficients) gives $L_X(\alpha \wedge \beta) = L_X\alpha \wedge \beta + \alpha \wedge L_X\beta$. Since $\Phi_t^*d = d\Phi_t^*$ (naturality of $d$,
Theorem~\ref{differential-forms:theorem-exterior-derivative}) and derivatives in $t$ and in the coordinates commute (the coefficients are smooth in $(t, p)$), $L_Xd = dL_X$. It is local,
as noted above.

\emph{Step 2: $D = d\iota_X + \iota_Xd$ is a local derivation of degree $0$ commuting with $d$.} For a $k$-form $\alpha$, using the antiderivation rules for $d$ and for $\iota_X$,
\[
\begin{aligned}
d\iota_X(\alpha \wedge \beta) &= d\iota_X\alpha \wedge \beta + (-1)^{k-1}\iota_X\alpha \wedge d\beta + (-1)^kd\alpha \wedge \iota_X\beta + \alpha \wedge d\iota_X\beta,\\
\iota_Xd(\alpha \wedge \beta) &= \iota_Xd\alpha \wedge \beta + (-1)^{k+1}d\alpha \wedge \iota_X\beta + (-1)^k\iota_X\alpha \wedge d\beta + \alpha \wedge \iota_Xd\beta,
\end{aligned}
\]
and adding, the middle terms cancel: $D(\alpha \wedge \beta) = D\alpha \wedge \beta + \alpha \wedge D\beta$. Moreover $dD = d\iota_Xd = Dd$, since $d \circ d = 0$.

\emph{Step 3: $L_X = D$.} On a function, $L_Xf = Xf = df(X) = \iota_Xdf = Df$, as $\iota_Xf = 0$. On an exact $1$-form $df$, both give $d(Xf)$, since both commute with $d$. Locally every form
is a sum of forms $f\,dg_1 \wedge \cdots \wedge dg_k$ (in a chart, $\omega_I\,dx^{i_1} \wedge \cdots \wedge dx^{i_k}$), and two local derivations of degree $0$ that agree on functions and on
exact $1$-forms agree on such products, hence everywhere.

\emph{Step 4: the derivative at time $t$.} By the group law, near a point and for small $s$, $\Phi_{t+s} = \Phi_s \circ \Phi_t$, so $\Phi_{t+s}^*\omega = \Phi_t^*(\Phi_s^*\omega)$, and differentiating at
$s = 0$ gives $\frac{d}{dt}\Phi_t^*\omega = \Phi_t^*L_X\omega$ ($\Phi_t^*$ is linear and acts pointwise, so it commutes with $\frac{d}{ds}$).

\emph{Step 5: brackets.} $L_XL_Y - L_YL_X$ is again a local derivation of degree $0$ commuting with $d$, and on functions it is $XY f - YXf = [X, Y]f = L_{[X,Y]}f$. As in Step 3 the
two derivations agree everywhere.
\end{proof}

By Step 4, the flow of $X$ preserves a form $\omega$, that is $\Phi_t^*\omega = \omega$ for all $t$, if and only if $L_X\omega = 0$. For example, the flow of a vector field on $\RR^n$
preserves the volume form $dx^1 \wedge \cdots \wedge dx^n$ exactly when $\operatorname{div}X = \sum_i\partial_iX^i = 0$
(Exercise~\ref{differential-forms:exercise-lie-derivative-volume}); in dimension $2$ the Hamiltonian vector fields are such
(Exercise~\ref{differential-forms:exercise-hamiltonian-area}).

\begin{example}\label{differential-forms:example-forms-r3}
On $\RR^3$ the Euclidean metric and orientation identify $1$-forms and $2$-forms with vector fields, by $F = (F_1, F_2, F_3) \mapsto F_1\,dx + F_2\,dy + F_3\,dz$ and
$F \mapsto F_1\,dy \wedge dz + F_2\,dz \wedge dx + F_3\,dx \wedge dy$, and $3$-forms with functions. Under these identifications $d$ on $0$-, $1$- and $2$-forms is the gradient, the curl and the divergence, and $d \circ d = 0$ becomes
$\operatorname{curl}\operatorname{grad} = 0$ and $\operatorname{div}\operatorname{curl} = 0$. Indeed $df = f_x\,dx + f_y\,dy + f_z\,dz$; for a $1$-form,
\[
d(F_1\,dx + F_2\,dy + F_3\,dz) = (\partial_yF_3 - \partial_zF_2)\,dy \wedge dz + (\partial_zF_1 - \partial_xF_3)\,dz \wedge dx + (\partial_xF_2 - \partial_yF_1)\,dx \wedge dy,
\]
since, for instance, $dF_1 \wedge dx = \partial_yF_1\,dy \wedge dx + \partial_zF_1\,dz \wedge dx$; and for a $2$-form,
$d(F_1\,dy \wedge dz + F_2\,dz \wedge dx + F_3\,dx \wedge dy) = (\partial_xF_1 + \partial_yF_2 + \partial_zF_3)\,dx \wedge dy \wedge dz$, because only the derivative in the missing variable
survives in each term.
\end{example}

\begin{definition}\label{differential-forms:definition-closed-exact}
A form $\omega$ is \emph{closed} if $d\omega = 0$ and \emph{exact} if $\omega = d\eta$ for some form $\eta$. Exact forms are closed, since $d \circ d = 0$.
\end{definition}

\begin{definition}\label{differential-forms:definition-line-integral}
For a $1$-form $\alpha$ on $M$ and a piecewise smooth curve $\gamma \colon [a, b] \to M$, the \emph{line integral} is $\int_\gamma\alpha = \int_a^b\alpha_{\gamma(t)}(\dot\gamma(t))\,dt$, the sum of the integrals of $\gamma^*\alpha$ over the smooth pieces.
\end{definition}

\begin{proposition}\label{differential-forms:proposition-line-integrals}
\begin{enumerate}
\item The line integral is unchanged under orientation-preserving reparametrisation $\gamma \circ \phi$, $\phi \colon [c, d] \to [a, b]$ increasing and piecewise smooth with $\phi(c) = a$, $\phi(d) = b$, and changes sign under reversal.
\item (Fundamental theorem for line integrals) $\int_\gamma df = f(\gamma(b)) - f(\gamma(a))$.
\item On a connected manifold, a $1$-form $\alpha$ is exact if and only if $\int_\gamma\alpha = 0$ for every closed piecewise smooth curve $\gamma$.
\end{enumerate}
\end{proposition}

\begin{proof}
(1) is the substitution rule (Corollary~\ref{real-analysis:corollary-integration-rules}) on each piece, as $(\gamma \circ \phi)^*\alpha = \phi^*\gamma^*\alpha$. (2) $\gamma^*df = d(f \circ \gamma) = (f \circ \gamma)'\,dt$, and
Theorem~\ref{real-analysis:theorem-ftc} applies on each piece; the intermediate values cancel. (3) If $\alpha = df$, closed curves give $0$ by (2). Conversely, fix $p_0$ and define $f(p) = \int_\gamma\alpha$ for any piecewise smooth curve
$\gamma$ from $p_0$ to $p$; such curves exist since $M$ is connected and locally path-connected by piecewise smooth curves (use straight segments in charts), and the value does not depend on $\gamma$, because for two
curves $\gamma$, $\gamma'$ the concatenation of $\gamma$ with the reverse of $\gamma'$ is closed. In a chart around $p$, with $p$ at the origin, computing $f$ along $\gamma$ followed by the segment from $p$ to $p + se_i$ gives
$f(p + se_i) - f(p) = \int_0^s\alpha_i(p + ue_i)\,du$, so $\partial_if = \alpha_i$ by Theorem~\ref{real-analysis:theorem-ftc}; hence $f$ is smooth and $df = \alpha$.
\end{proof}

\begin{counterexample}\label{differential-forms:counterexample-closed-not-exact}
Closed forms need not be exact. On $M = \RR^2 \setminus \{0\}$, the \emph{angle form} $\alpha = (x\,dy - y\,dx)/(x^2 + y^2)$ is closed, as a direct computation shows, but along the unit circle $\gamma(t) = (\cos t, \sin t)$,
$\gamma^*\alpha = (\cos^2t + \sin^2t)\,dt = dt$, so $\int_\gamma\alpha = 2\pi \neq 0$, and $\alpha$ is not exact by Proposition~\ref{differential-forms:proposition-line-integrals}(3). On the half-plane $\{x > 0\}$ it is exact, equal to
$d\arctan(y/x)$: exactness is a global question, measured by de Rham cohomology (Chapter~\ref{de-rham}).
\end{counterexample}

\section{Orientations}\label{differential-forms:section-orientations}

\begin{definition}\label{differential-forms:definition-orientation}
An \emph{orientation} of a smooth $n$-manifold $M$ is an equivalence class of nowhere vanishing $n$-forms, two being equivalent if they differ by a positive function; $M$ is \emph{orientable} if such a
form exists, and \emph{oriented} once one is chosen. A chart is \emph{positively oriented} if $dx^1 \wedge \cdots \wedge dx^n$ represents the orientation.
\end{definition}

\begin{proposition}\label{differential-forms:proposition-orientation}
$M$ is orientable if and only if it has an atlas whose transition maps have positive Jacobian determinant, if and only if the line bundle $\det T^*M$ is trivial. For connected $M$ there are then exactly two
orientations. This agrees with the topological notion of Definition~\ref{poincare-duality:definition-orientation-covering}.
\end{proposition}

\begin{proof}
A nowhere vanishing $n$-form $\mu$ gives such an atlas: take charts with $\mu = f\,dx^1 \wedge \cdots \wedge dx^n$, $f > 0$ (compose with a reflection if necessary), and on overlaps the transition Jacobian is
the ratio of two positive functions. Conversely, given such an atlas, patch the forms $dx^1 \wedge \cdots \wedge dx^n$ with a partition of unity: the sum is nowhere vanishing because in each chart it is a
positive combination of positive multiples of $dx^1 \wedge \cdots \wedge dx^n$. A nowhere vanishing section of the line bundle $\det T^*M = \Lambda^nT^*M$ is the same as a trivialisation. For the last
statement, in a positively oriented chart a generator of $H_n(M \mid p; \ZZ)$ is determined by the orientation of $\RR^n$, and this choice is locally constant, giving a section of the orientation covering;
conversely such a section determines a compatible atlas.
\end{proof}

\begin{counterexample}\label{differential-forms:counterexample-mobius}
The open M\"obius band is not orientable. Let $\ZZ$ act on $\RR \times (-1, 1)$ by $g^k(x, t) = (x + k, (-1)^kt)$, freely and properly, and let $B$ be the smooth quotient (Theorem~\ref{smooth-manifolds:theorem-smooth-quotient}). A
nowhere vanishing $2$-form $\mu$ on $B$ would pull back to a nowhere vanishing form $f\,dx \wedge dt$ on $\RR \times (-1, 1)$ invariant under $g$. As $g^*(f\,dx \wedge dt) = (f \circ g)\,dx \wedge d(-t) = -(f \circ g)\,dx \wedge dt$,
invariance means $f \circ g = -f$. But $f$ is continuous and nowhere zero on the connected set $\RR \times (-1, 1)$, hence of constant sign (Corollary~\ref{general-topology:corollary-ivt}), which contradicts $f(g(x, t)) = -f(x, t)$.
\end{counterexample}

For surfaces in $\RR^3$ one usually thinks of an orientation as a choice of side, that is, of a unit normal vector varying continuously along the surface. This is
the same notion.

\begin{proposition}\label{differential-forms:proposition-orientation-normal}
Let $S \subseteq \RR^{n+1}$ be an $n$-dimensional submanifold. A \emph{unit normal field} is a continuous map $\nu \colon S \to \RR^{n+1}$ with $|\nu(p)| = 1$ and
$\nu(p) \perp T_pS$ for all $p$. Unit normal fields correspond bijectively to orientations of $S$: to $\nu$ corresponds the orientation of the nowhere vanishing
$n$-form $\mu_p(v_1, \ldots, v_n) = \det(\nu(p), v_1, \ldots, v_n)$, for which a basis $(v_1, \ldots, v_n)$ of $T_pS$ is positive exactly when $(\nu(p), v_1, \ldots, v_n)$ is a
positive basis of $\RR^{n+1}$. In particular $S$ is orientable if and only if it has a unit normal field, and every unit normal field is smooth. For a surface in $\RR^3$, a
basis $(v, w)$ of $T_pS$ is positive exactly when $v \times w$ is a positive multiple of $\nu(p)$.
\end{proposition}

\begin{proof}
Near a point of $S$ write $S$ as a regular level set $F^{-1}(0)$ with $F$ real valued (Proposition~\ref{smooth-manifolds:proposition-submanifold-criteria}(4)). Then
$T_pS = \ker d_pF = \nabla F(p)^\perp$ (Theorem~\ref{smooth-manifolds:theorem-regular-value}), so the normal line at $p$ is spanned by $\nabla F(p) \neq 0$, and a unit normal field
is there $\varepsilon\,\nabla F/|\nabla F|$ with $\varepsilon = \pm1$ continuous, hence locally constant; so $\nu$ is smooth. Given $\nu$, $\mu$ is a smooth $n$-form on $S$ (in a
parametrisation its coefficient is $\det(\nu, x_1, \ldots, x_n)$), and it vanishes nowhere, because $\nu(p) \notin T_pS$, so $(\nu(p), v_1, \ldots, v_n)$ is a basis of $\RR^{n+1}$
whenever $(v_i)$ is one of $T_pS$. Conversely, given an orientation, exactly one of the two unit normals $\pm\nu(p)$ makes $\det(\nu(p), v_1, \ldots, v_n) > 0$ for positive bases
$(v_i)$, the choice not depending on the positive basis, as two positive bases differ by a matrix of positive determinant. Near a point, in the notation above,
$\nu = \varepsilon\nabla F/|\nabla F|$ with $\varepsilon$ the sign of $\det(\nabla F, \partial_1, \ldots, \partial_n)$ for the coordinate fields of a positive chart, a continuous nowhere zero
function; so $\nu$ is continuous. The two constructions are inverse to each other. For $n = 2$, $\det(\nu, v, w) = \langle\nu, v \times w\rangle$ and $v \times w$ is normal to
$T_pS$.
\end{proof}

\begin{remark}\label{differential-forms:remark-moving-normal}
So a hypersurface is orientable when a unit normal can be carried continuously around every closed curve and comes back to itself. The regular level sets
$F^{-1}(c) \subseteq \RR^{n+1}$ are orientable, with $\nu = \nabla F/|\nabla F|$; for instance spheres, with the outward normal. For the M\"obius band embedded in
$\RR^3$ as the image of $x(s, t) = ((1 + t\cos\frac s2)\cos s, (1 + t\cos\frac s2)\sin s, t\sin\frac s2)$, $|t| < \frac12$, the normal along the central circle $t = 0$ is
$x_s \times x_t = (\cos s\sin\frac s2, \sin s\sin\frac s2, -\cos\frac s2)$, which is $(0, 0, -1)$ at $s = 0$ and $(0, 0, 1)$ at $s = 2\pi$, at the same point $(1, 0, 0)$: going once
around reverses the normal, so there is no unit normal field, in accordance with Counterexample~\ref{differential-forms:counterexample-mobius}. The induced orientation of a
boundary (Definition~\ref{differential-forms:definition-manifold-with-boundary} below) is the same idea with the outward normal: for a domain in $\RR^3$ bounded by a surface, the boundary
is oriented by the outward unit normal, which is the convention of the divergence theorem.
\end{remark}

\section{Integration and Stokes' theorem}\label{differential-forms:section-integration}

Let $\tau \colon V \to U$ be a diffeomorphism between open subsets of $\RR^n$, with coordinates $y$ on $V$ and $x$ on $U$. For a function $f$ on $U$, the change of variables formula
(Theorem~\ref{measure-theory:theorem-change-of-variables}) says $\int_Uf\,d\lambda = \int_V(f \circ \tau)\,|\det D\tau|\,d\lambda$, while the pullback of the $n$-form $f\,dx^1 \wedge \cdots \wedge dx^n$
is $(f \circ \tau)\det D\tau\,dy^1 \wedge \cdots \wedge dy^n$ (Proposition~\ref{differential-forms:proposition-forms-descriptions}). So if we integrate an $n$-form by integrating its coefficient,
the answer does not change under orientation-preserving changes of coordinates ($\det D\tau > 0$) and changes sign under orientation-reversing ones. This is the whole reason for
the following definition.

\begin{definition}\label{differential-forms:definition-integral}
Let $M$ be an oriented $n$-manifold and $\omega$ an $n$-form with compact support contained in a positively oriented chart $\varphi \colon U \to \RR^n$. Write $(\varphi^{-1})^*\omega = f\,dx^1 \wedge \cdots \wedge dx^n$
and set $\int_M\omega = \int_{\varphi(U)}f\,d\lambda$. For general compactly supported $\omega$, choose a partition of unity $(\psi_i)$ subordinate to a cover by positively oriented charts and set
$\int_M\omega = \sum_i\int_M\psi_i\omega$.
\end{definition}

\begin{proposition}\label{differential-forms:proposition-integral-well-defined}
The integral is well defined, linear, and satisfies $\int_M\omega = \int_N f^*\omega$ for an orientation-preserving diffeomorphism $f \colon N \to M$; it changes sign if the orientation is reversed.
\end{proposition}

\begin{proof}
\emph{One chart.} Let $\omega$ have compact support in the domains of two positively oriented charts $\varphi$ and $\psi$, and $\tau = \varphi \circ \psi^{-1}$, a diffeomorphism between open subsets of
$\RR^n$ with $\det D\tau > 0$. If $(\varphi^{-1})^*\omega = f\,dx$, with $dx = dx^1 \wedge \cdots \wedge dx^n$, then $(\psi^{-1})^*\omega = \tau^*(\varphi^{-1})^*\omega = (f \circ \tau)\det D\tau\,dy$, and by the
change of variables formula $\int(f \circ \tau)\det D\tau\,d\lambda = \int(f \circ \tau)|\det D\tau|\,d\lambda = \int f\,d\lambda$. (The supports lie in the common domain, and the integrands vanish
outside it.)

\emph{Partitions of unity.} The sum in the definition is finite, because the compact support of $\omega$ meets only finitely many supports of a locally finite partition of unity, and
each term is defined by the first case. For two partitions $(\psi_i)$ and $(\chi_j)$, each $\psi_i\chi_j\omega$ is supported in the domains of the charts of both, so
$\sum_i\int\psi_i\omega = \sum_{i,j}\int\chi_j\psi_i\omega = \sum_j\int\chi_j\omega$, using the linearity of the one-chart integral.

\emph{Properties.} Linearity follows from that of the one-chart integral. If $f \colon N \to M$ is an orientation-preserving diffeomorphism, compositions $\varphi \circ f$ of positive charts of
$M$ with $f$ are positive charts of $N$, and $\psi_i \circ f$ is a partition of unity; the coefficient of $f^*\omega$ in the chart $\varphi \circ f$ is that of $\omega$ in $\varphi$. So the defining sums
agree. Reversing the orientation replaces each positive chart $\varphi$ by $r \circ \varphi$, with the reflection $r(x^1, \ldots, x^n) = (-x^1, x^2, \ldots, x^n)$, which changes the sign of every
one-chart integral.
\end{proof}

In practice one does not use partitions of unity: one parametrises $M$ up to a set of measure zero.

\begin{proposition}\label{differential-forms:proposition-integral-parametrisation}
Let $M$ be an oriented $n$-manifold, $U \subseteq \RR^n$ open and $x \colon U \to M$ an orientation-preserving diffeomorphism onto an open subset whose complement $M \setminus x(U)$ has
measure zero, that is, its image in every chart has Lebesgue measure $0$. For a compactly supported $n$-form $\omega$ write $x^*\omega = f\,du^1 \wedge \cdots \wedge du^n$. Then $f$ is integrable
and
\[
\int_M\omega = \int_Uf\,d\lambda .
\]
\end{proposition}

\begin{proof}
Choose a partition of unity $(\psi_i)$ subordinate to positively oriented charts $(U_i, \varphi_i)$; finitely many $\psi_i$ meet the support of $\omega$. In the one-chart integral of $\psi_i\omega$ we may
remove the null set $\varphi_i(U_i \setminus x(U))$. On the rest, $\tau_i = \varphi_i \circ x$ is an orientation-preserving diffeomorphism from $x^{-1}(U_i)$ onto $\varphi_i(U_i \cap x(U))$, so by the
change of variables formula, as in the proof of Proposition~\ref{differential-forms:proposition-integral-well-defined}, $\int_M\psi_i\omega = \int_{x^{-1}(U_i)}f_i\,d\lambda$, where
$x^*(\psi_i\omega) = f_i\,du$. Summing over the finitely many $i$, and using $\sum_if_i = f$ on $U$, gives the claim; each $f_i$ is integrable, so $f$ is.
\end{proof}

\begin{example}\label{differential-forms:example-integrals}
\begin{enumerate}
\item \emph{Curves.} If $\gamma \colon (a, b) \to M$ is an embedding onto a curve $C$, oriented by $\gamma$, then $\int_C\alpha = \int_a^b\alpha(\dot\gamma(t))\,dt$ is the line integral of
Definition~\ref{differential-forms:definition-line-integral}.
\item \emph{The sphere.} Orient $S^2$ by its outward normal (Proposition~\ref{differential-forms:proposition-orientation-normal}); the spherical coordinates
$x(\vartheta, \varphi)$ on $(0, \pi) \times (0, 2\pi)$ are positively oriented, since $x_\vartheta \times x_\varphi = \sin\vartheta\,x$ points outward, and they miss only a closed half great
circle, which has measure zero (it is a segment in the chart given by $\varphi \in (-\pi, \pi)$, together with two points). By
Example~\ref{differential-forms:example-form-computations}(3) and Proposition~\ref{differential-forms:proposition-integral-parametrisation},
\[
\int_{S^2}x\,dy \wedge dz + y\,dz \wedge dx + z\,dx \wedge dy = \int_0^{2\pi}\int_0^\pi\sin\vartheta\,d\vartheta\,d\varphi = 4\pi .
\]
With the inner normal the same integral is $-4\pi$.
\item \emph{The angle form.} $\int_{S^1}\alpha = 2\pi$ for $\alpha = (x\,dy - y\,dx)/(x^2 + y^2)$ and the counterclockwise orientation of $S^1 \subseteq \RR^2$
(Counterexample~\ref{differential-forms:counterexample-closed-not-exact}).
\end{enumerate}
\end{example}

\begin{definition}\label{differential-forms:definition-manifold-with-boundary}
A \emph{smooth $n$-manifold with boundary} is defined as in Definition~\ref{smooth-manifolds:definition-manifold} with charts taking values in open subsets of the half-space
$\mathbb{H}^n = \{x \in \RR^n : x^n \geq 0\}$; the \emph{boundary} $\partial M$ is the set of points corresponding to $x^n = 0$, a smooth $(n-1)$-manifold. If $M$ is oriented, $\partial M$ carries the
\emph{induced orientation}, defined by declaring a chart of $\partial M$ positively oriented when the $n$-form $dx^1 \wedge \cdots \wedge dx^n$ of a positively oriented chart of $M$ is $(-1)^n\,\nu \wedge
(dx^1 \wedge \cdots \wedge dx^{n-1})$ for an outward-pointing covector $\nu$; concretely, $(x^1, \ldots, x^{n-1})$ is positively oriented on $\partial M$ when $(-1)^n(x^1, \ldots, x^n)$ is on $M$.
\end{definition}

Here a map defined on an open subset of $\mathbb{H}^n$ (open for the subspace topology) is \emph{smooth} if near each point it is the restriction of a smooth map defined on an open subset of $\RR^n$;
the transition maps between charts are required to be smooth in this sense. Tangent spaces, vector fields and forms are defined as for manifolds; at a boundary point the tangent space is
still $n$-dimensional, and a tangent vector is \emph{outward-pointing} if in a chart its $n$-th component is negative. The rule for the induced orientation says: a basis
$(v_1, \ldots, v_{n-1})$ of $T_p\partial M$ is positive if and only if $(w, v_1, \ldots, v_{n-1})$ is a positive basis of $T_pM$ for an outward vector $w$ (\emph{outward vector first}).
Indeed, in a positive chart $w = -\partial_n$ is outward, and $\det(-e_n, e_1, \ldots, e_{n-1}) = (-1)^n$.

\begin{lemma}\label{differential-forms:lemma-boundary-well-defined}
Whether a point of a manifold with boundary lies on the boundary does not depend on the chart. The restrictions of the charts to $\partial M$ form an atlas of an $(n-1)$-manifold without
boundary, and $M \setminus \partial M$ is an $n$-manifold without boundary.
\end{lemma}

\begin{proof}
Suppose a chart $\varphi$ sends $p$ to a point with $x^n = 0$ and a chart $\psi$ sends it to a point with $x^n > 0$. Shrinking the domain, $\psi$ maps onto an open subset $W$ of $\RR^n$ contained in
$\{x^n > 0\}$, and $\tau = \varphi \circ \psi^{-1} \colon W \to \mathbb{H}^n$ is smooth with a smooth inverse, which near $\varphi(p)$ extends to a smooth map $\sigma$ on an open subset of $\RR^n$.
From $\sigma \circ \tau = \id$ near $\psi(p)$, the chain rule gives $D\sigma(\varphi(p))D\tau(\psi(p)) = I$, so $D\tau(\psi(p))$ is invertible. By the inverse function theorem
(Theorem~\ref{real-analysis:theorem-inverse-function}), $\tau$ maps a neighbourhood of $\psi(p)$ onto an open subset of $\RR^n$ containing $\varphi(p)$; but every neighbourhood of $\varphi(p)$ in
$\RR^n$ contains points with $x^n < 0$, which are not in $\mathbb{H}^n$. This contradiction proves the first statement. The restrictions of the transition maps to $\{x^n = 0\} = \RR^{n-1}$ are
smooth with smooth inverses, which gives the atlas of $\partial M$, and the charts restricted to the interior points take values in the open set $\{x^n > 0\}$ of $\RR^n$.
\end{proof}

\begin{example}\label{differential-forms:example-boundaries}
\begin{enumerate}
\item $M = [a, b]$, with the orientation of $\RR$. At $b$ the outward direction is $+\partial_t$ and at $a$ it is $-\partial_t$; a point is oriented by a sign, and ``outward vector first''
gives $\partial[a, b] = \{b\} - \{a\}$: the sign $+1$ at $b$ and $-1$ at $a$.
\item The closed unit disc $\bar D \subseteq \RR^2$: at a boundary point $x$ the vector $x$ points outward, and $(x, v)$ is a positive basis of $\RR^2$ exactly when $v$ is a positive multiple of
$(-x_2, x_1)$. So the boundary circle is oriented counterclockwise: the disc is on the left.
\item The closed ball in $\RR^3$: its boundary $S^2$ is oriented by the outward normal, in the sense of Proposition~\ref{differential-forms:proposition-orientation-normal}.
\item The half-space $\mathbb{H}^n$ itself: $\partial\mathbb{H}^n = \RR^{n-1}$ with $(-1)^n$ times its standard orientation; for $n = 2$ the boundary of the upper half-plane is the
real line oriented from left to right.
\end{enumerate}
\end{example}

\begin{theorem}[Stokes]\label{differential-forms:theorem-stokes}
Let $M$ be an oriented $n$-manifold with boundary and $\omega$ an $(n-1)$-form with compact support. Then
\[
\int_Md\omega = \int_{\partial M}\omega,
\]
where $\partial M$ has the induced orientation. If $\partial M = \emptyset$, the integral of an exact compactly supported form vanishes.
\end{theorem}

\begin{proof}
\emph{Reduction to one chart.} Let $(\psi_i)$ be a partition of unity subordinate to positively oriented charts. Only finitely many $\psi_i\omega$ are nonzero, $\omega = \sum_i\psi_i\omega$ and
$d\omega = \sum_id(\psi_i\omega)$. Both sides of the formula are linear in $\omega$, so it suffices to prove it for each $\psi_i\omega$, that is, for $\omega$ supported in a single positively oriented
chart. Via the chart, $M$ becomes an open subset of $\RR^n$ or of $\mathbb{H}^n$, and extending $\omega$ by $0$, we may assume $M = \RR^n$ or $M = \mathbb{H}^n$ and $\omega$ has compact support.

\emph{The computation.} Write
$\omega = \sum_i(-1)^{i-1}f_i\,dx^1 \wedge \cdots \wedge \widehat{dx^i} \wedge \cdots \wedge dx^n$; then $d\omega = (\sum_i\partial_if_i)\,dx^1 \wedge \cdots \wedge dx^n$, because
$d f_i \wedge dx^1 \wedge \cdots \wedge \widehat{dx^i} \wedge \cdots \wedge dx^n = (-1)^{i-1}\partial_if_i\,dx^1 \wedge \cdots \wedge dx^n$. By Fubini's theorem (Theorem~\ref{measure-theory:theorem-fubini})
we may integrate $\partial_if_i$ first in the variable $x^i$. For $i < n$, and also for $i = n$ when $M = \RR^n$, the integral over the whole line vanishes by the fundamental theorem of calculus,
as $f_i$ has compact support. So $\int_{\RR^n}d\omega = 0 = \int_{\partial\RR^n}\omega$. For $M = \mathbb{H}^n$ the only surviving term is
\[
\int_{\mathbb{H}^n}\partial_nf_n\,d\lambda = \int_{\RR^{n-1}}\left(\int_0^\infty\partial_nf_n\,dx^n\right)dx^1 \cdots dx^{n-1} = -\int_{\RR^{n-1}}f_n(x', 0)\,dx' .
\]
On $\partial\mathbb{H}^n = \{x^n = 0\}$ the forms containing $dx^n$ restrict to $0$, so $\omega|_{\partial\mathbb{H}^n} = (-1)^{n-1}f_n\,dx^1 \wedge \cdots \wedge dx^{n-1}$, and the induced orientation is
$(-1)^n$ times the standard one (Example~\ref{differential-forms:example-boundaries}(4)). Hence $\int_{\partial\mathbb{H}^n}\omega = (-1)^n(-1)^{n-1}\int f_n(x', 0)\,dx' = -\int f_n(x', 0)\,dx'$,
which is the left side.
\end{proof}

\begin{counterexample}\label{differential-forms:counterexample-stokes-support}
Stokes' theorem needs compact support. On $M = \RR$, without boundary, the $0$-form $\arctan$ has $\int_\RR d\arctan = \int_\RR\frac{dx}{1 + x^2} = \pi \neq 0$. On $M = [0, +\infty)$, with boundary the point $0$ carrying the induced orientation $-1$,
the constant function $1$ has $\int_Md1 = 0$ while $\int_{\partial M}1 = -1$.
\end{counterexample}

\begin{example}\label{differential-forms:example-stokes-dimension-one}
For $M = [a, b]$ with its standard orientation, $\partial M = \{b\} - \{a\}$ with the induced orientation, and Stokes' theorem for a function $f$ is the fundamental theorem of calculus, $\int_a^bf' = f(b) - f(a)$. For a compact region
$M \subseteq \RR^2$ with smooth boundary and $\omega = P\,dx + Q\,dy$, it is Green's formula $\int_M(\partial_xQ - \partial_yP)\,dx\,dy = \int_{\partial M}P\,dx + Q\,dy$, the boundary being traversed with $M$ on the left.
\end{example}

\begin{example}\label{differential-forms:example-classical}
For $M \subseteq \RR^3$ a compact region with smooth boundary, Stokes' theorem specialises to the divergence theorem $\int_M\operatorname{div}F\,dV = \int_{\partial M}\langle F, \nu\rangle\,dA$ and to the classical
Stokes and Green formulas, after identifying vector fields with $1$- or $2$-forms using the Euclidean metric and orientation (Theorem~\ref{multilinear-algebra:theorem-hodge-star}). On a compact oriented
manifold without boundary, $\int_M d\omega = 0$, so integration descends to de Rham cohomology (Chapter~\ref{de-rham}).
\end{example}

\begin{example}\label{differential-forms:example-stokes-checks}
\begin{enumerate}
\item \emph{Area by a line integral.} For the unit disc $\bar D$ and $\omega = x\,dy$, $d\omega = dx \wedge dy$, so $\int_{\partial\bar D}x\,dy = \int_{\bar D}dx \wedge dy = \pi$. Directly, with
$x = \cos t$, $y = \sin t$: $\int_0^{2\pi}\cos^2t\,dt = \pi$. More generally the area of a plane region is $\frac12\int_{\partial R}(x\,dy - y\,dx)$.
\item \emph{Volume by a flux.} For the unit ball $B$ and $\omega = x\,dy \wedge dz + y\,dz \wedge dx + z\,dx \wedge dy$, $d\omega = 3\,dx \wedge dy \wedge dz$, and Stokes' theorem gives
$\int_{S^2}\omega = 3\operatorname{vol}(B) = 4\pi$, which is the value computed directly in Example~\ref{differential-forms:example-integrals}(2). So $\operatorname{vol}(B) = 4\pi/3$.
\item \emph{A closed form on an annulus.} For the annulus $A = \{1 \leq |x| \leq 2\}$ and the angle form $\alpha$, $d\alpha = 0$ and $\partial A$ is the outer circle counterclockwise and the
inner circle clockwise, so $\int_{\partial A}\alpha = 2\pi - 2\pi = 0$, as it must be. The punctured disc is not compact, and there Stokes' theorem does not apply: the closed form $\alpha$
has integral $2\pi$ over the boundary circle.
\end{enumerate}
\end{example}

\begin{corollary}\label{differential-forms:corollary-volume}
On a compact oriented $n$-manifold $M$ without boundary, $\int_M \colon \Omega^n(M) \to \RR$ is surjective, and $\int_M\mu > 0$ for any $\mu$ representing the orientation. Hence $M$ carries a \emph{volume
form}, and any two differ by a positive function.
\end{corollary}

\begin{proof}
In a positively oriented chart the integral of a nonnegative bump $n$-form is positive; a form representing the orientation is positive in every positively oriented chart, so its integral is positive by the
definition via a partition of unity.
\end{proof}

\section{Exercises}\label{differential-forms:section-exercises}

\begin{exercise}\label{differential-forms:exercise-pullback-computation}
Compute the pullback of $dx \wedge dy$ under polar coordinates $(r, \theta) \mapsto (r\cos\theta, r\sin\theta)$, and the exterior derivative of $\alpha = (x\,dy - y\,dx)/(x^2 + y^2)$ on $\RR^2 \setminus \{0\}$.
Show that $\alpha$ is closed but not exact, by integrating over the unit circle.
\end{exercise}

\begin{solution}
The pullback is $r\,dr \wedge d\theta$. A computation gives $d\alpha = 0$, while $\int_{S^1}\alpha = 2\pi \neq 0$; an exact form would integrate to $0$ over the closed curve $S^1$ by Stokes' theorem.
\end{solution}

\begin{exercise}\label{differential-forms:exercise-orientability-projective}
Show that $\RR\PP^n$ is orientable if and only if $n$ is odd, by computing the effect of the antipodal map on a volume form of $S^n$ (compare
Exercise~\ref{poincare-duality:exercise-projective-orientability}).
\end{exercise}

\begin{exercise}\label{differential-forms:exercise-lie-derivative-volume}
Let $\mu$ be a volume form and $X$ a vector field. Show that $L_X\mu = (\operatorname{div}X)\mu$ defines the divergence, that $\operatorname{div}X = \sum_i\partial_iX^i$ in coordinates with
$\mu = dx^1 \wedge \cdots \wedge dx^n$, and that the flow of $X$ preserves $\mu$ if and only if $\operatorname{div}X = 0$.
\end{exercise}

\begin{solution}
$L_X\mu$ is an $n$-form, hence a function times $\mu$. By Cartan's formula $L_X\mu = d\iota_X\mu$, and computing $\iota_X\mu$ in coordinates gives $\sum_i(-1)^{i-1}X^i\,dx^1 \wedge \cdots \widehat{dx^i} \cdots \wedge dx^n$,
whose differential is $(\sum_i\partial_iX^i)\mu$. The flow preserves $\mu$ if and only if $\frac{d}{dt}\Phi_t^*\mu = \Phi_t^*L_X\mu$ vanishes.
\end{solution}

\begin{exercise}\label{differential-forms:exercise-stokes-degree}
Let $f \colon M \to N$ be a smooth map between compact oriented $n$-manifolds without boundary. Show that $\int_Mf^*\omega = \deg(f)\int_N\omega$ for a suitable integer $\deg(f)$ independent of $\omega$, using
Sard's theorem and the local form of $f$ near regular values.
\end{exercise}

\begin{exercise}\label{differential-forms:exercise-exact-primitive}
On $\RR^3$ let $\omega = yz\,dx + zx\,dy + xy\,dz$. Show that $d\omega = 0$ and find $f$ with $df = \omega$. Compute also $d(e^{xy}\,dx \wedge dz)$.
\end{exercise}

\begin{solution}
$d\omega = (z\,dy + y\,dz) \wedge dx + (z\,dx + x\,dz) \wedge dy + (y\,dx + x\,dy) \wedge dz$, and the six terms cancel in pairs, for example $z\,dy \wedge dx + z\,dx \wedge dy = 0$. A primitive is
$f = xyz$. Next, $d(e^{xy}dx \wedge dz) = (ye^{xy}dx + xe^{xy}dy) \wedge dx \wedge dz = xe^{xy}\,dy \wedge dx \wedge dz = -xe^{xy}\,dx \wedge dy \wedge dz$.
\end{solution}

\begin{exercise}\label{differential-forms:exercise-spherical-volume}
Show that the pullback of $dx \wedge dy \wedge dz$ under spherical coordinates $(r, \vartheta, \varphi) \mapsto (r\sin\vartheta\cos\varphi, r\sin\vartheta\sin\varphi, r\cos\vartheta)$ is
$r^2\sin\vartheta\,dr \wedge d\vartheta \wedge d\varphi$, and compute the volume of the ball of radius $R$ as the integral of $dx \wedge dy \wedge dz$.
\end{exercise}

\begin{solution}
The pullback is the Jacobian determinant times $dr \wedge d\vartheta \wedge d\varphi$. Expanding the $3 \times 3$ determinant along the row of $z = r\cos\vartheta$, whose entries are
$\cos\vartheta$, $-r\sin\vartheta$, $0$, gives $\cos\vartheta \cdot r^2\sin\vartheta\cos\vartheta + r\sin\vartheta \cdot r\sin^2\vartheta = r^2\sin\vartheta$. By
Proposition~\ref{differential-forms:proposition-integral-parametrisation}, the volume is $\int_0^{2\pi}\int_0^\pi\int_0^Rr^2\sin\vartheta\,dr\,d\vartheta\,d\varphi = \frac43\pi R^3$.
\end{solution}

\begin{exercise}\label{differential-forms:exercise-stokes-hemisphere}
Verify Stokes' theorem for $\omega = x\,dy - y\,dx$ on the closed upper hemisphere $H = S^2 \cap \{z \geq 0\}$, oriented by the outward normal, by computing both sides.
\end{exercise}

\begin{solution}
$d\omega = 2\,dx \wedge dy$. The parametrisation $(u, v) \mapsto (u, v, \sqrt{1 - u^2 - v^2})$ of the open upper hemisphere by the open unit disc is positively oriented, as its normal
$x_u \times x_v$ has positive third component, and the pullback of $dx \wedge dy$ is $du \wedge dv$; so $\int_Hd\omega = 2\pi$
(Proposition~\ref{differential-forms:proposition-integral-parametrisation}; the equator has measure zero in $H$). The boundary is the equator. At $(1, 0, 0)$ an outward vector of $H$ is $-e_3$,
and $(-e_3, e_2)$ is positive for the outward normal $e_1$, since $\det(e_1, -e_3, e_2) = 1$; so the equator is traversed counterclockwise seen from above, and
$\int_{\partial H}x\,dy - y\,dx = \int_0^{2\pi}(\cos^2t + \sin^2t)\,dt = 2\pi$.
\end{solution}

\begin{exercise}\label{differential-forms:exercise-lie-derivative-rotation}
On $\RR^3$ let $X = x\,\partial_y - y\,\partial_x$ be the infinitesimal rotation about the $z$-axis. Compute $L_X$ of $dx$, $dy$, $dz$, $x\,dy - y\,dx$ and $dx \wedge dy$, and interpret the results.
\end{exercise}

\begin{solution}
$L_Xdx = d(Xx) = d(-y) = -dy$, $L_Xdy = d(Xy) = dx$, $L_Xdz = 0$. Then $L_X(x\,dy) = (Xx)\,dy + x\,L_Xdy = -y\,dy + x\,dx$ and $L_X(y\,dx) = x\,dx - y\,dy$, so
$L_X(x\,dy - y\,dx) = 0$, and $L_X(dx \wedge dy) = -dy \wedge dy + dx \wedge dx = 0$. The rotations preserve the angular form and the area form of the planes orthogonal to the axis, and the
height $z$.
\end{solution}

\begin{exercise}\label{differential-forms:exercise-hamiltonian-area}
For a smooth function $H$ on $\RR^2$ let $X_H = \partial_yH\,\partial_x - \partial_xH\,\partial_y$. Show that $\iota_{X_H}(dx \wedge dy) = dH$, that $X_HH = 0$, and deduce that the flow of $X_H$
preserves areas and the level sets of $H$.
\end{exercise}

\begin{solution}
$\iota_{X_H}(dx \wedge dy) = \partial_yH\,dy + \partial_xH\,dx = dH$. By Cartan's formula $L_{X_H}(dx \wedge dy) = d\iota_{X_H}(dx \wedge dy) = d\,dH = 0$, so the flow preserves $dx \wedge dy$
(Proposition~\ref{differential-forms:proposition-cartan}). And $X_HH = \partial_yH\,\partial_xH - \partial_xH\,\partial_yH = 0$, so $H$ is constant along the flow lines. This is the
two-dimensional case of Liouville's theorem in Hamiltonian mechanics.
\end{solution}
